\chapter{Conclusion}
\label{ch:conclusion}

This project built a teleoperation system for two seven-degree-of-freedom arms
mounted on a worn backpack frame, commanded from an instrumented mannequin
master, with a safety architecture appropriate to a manipulator standing
beside a person and a shared-autonomy layer above it.

\section{Against the objectives}

\emph{A master, built and characterised.} Fourteen potentiometer channels, two
inertial sensors, two force sensors and two buttons on a single
microcontroller, emitting a text frame at \SI{100}{\hertz}. Characterisation
found six of fourteen channels usable and reversed two long-standing verdicts,
both of which had come from computing statistics across transmitted rows
rather than distinct sensor updates.

\emph{A mapping that degrades predictably.} Position in spherical form with
each component taken from the sensor that observes it best. Vertical and
fore-and-aft motion track correctly on both arms. Lateral does not, and the
measurements show why no axis map or recalibration can fix it. The cost of
each channel reduction is quantified on \num{20430} recorded frames, and the
endpoint of the ladder reports position unavailable rather than inventing one.

\emph{A safety architecture validated by injection.} Fourteen of fourteen
injected faults handled. The two that took several attempts both failed for
the same underlying reason, which is that a mechanism was keyed on the wrong
observable: a dead-man on message arrival cannot detect a source that keeps
publishing stale data.

\emph{A workspace characterisation.} The two arms share no reachable
workspace under the as-built mount, which blocks inter-arm handover for
geometric reasons no software change can address. A buildable mount
modification that would raise the shared cell count from three to twenty-two
of sixty-three is quantified and deliberately not applied, because applying it
invalidates the home pose, the workspace anchor and every clearance figure
derived from them.

\emph{A verification methodology.} Ten repeats over densified paths rather
than single calls at endpoints, and an instrument validated against a known
answer before a surprising measurement taken with it is reported. The
five-task verification reproduced exactly across runs, which is the first
reachability figure in this project to do so.

\emph{Identification of what blocks progress.} Five items, listed in
\cref{ch:discussion}. The first is two soldered joints.

\section{What was not achieved}

No result in this thesis was validated on a physical Kinova arm. The task set
is verified geometrically and has never been performed. The two-person study
the architecture was designed around was not run, and could not have been:
the direct teleoperation baseline depends on channels that are currently
incoherent.

The planning report's haptic objective was not delivered and, as
\cref{ch:discussion} argues, is not deliverable on this platform through the
available control interface.

\section{Future work, in the order it should be done}

\begin{enumerate}
  \item \textbf{Repair left j2 and j4, and re-run the channel baseline after
    each attempt.} These two are named by the observability regression rather
    than chosen: with both frozen the left arm retains no reach observable at
    all, and repairing any other channel will not restore one. Re-running
    after each attempt rather than once at the end catches a repair that
    breaks a working channel.
  \item \textbf{Correct the filter seed in the firmware} so the potentiometers
    do not spend \SI{440}{\milli\second} settling from a railed value after
    every board reset. The change is one line and does not alter the filter's
    steady state.
  \item \textbf{Read the right arm's home joint angles from hardware.} They
    have never been read, and the workspace geometry depends on them.
  \item \textbf{Bring up a second inertial sensor per arm on the secondary
    bus.} The elbow angle then comes from the relative orientation of two
    segments with no potentiometer, and the shared yaw ambiguity cancels in a
    relative measurement, which is why two sensors can replace the elbow and
    one cannot. It costs two potentiometer channels and reaches
    \SI{63.8}{\milli\metre} on a single remaining pot.
  \item \textbf{Re-park the right arm and re-derive the mount.} This is the
    only route to a shared workspace, and therefore the only route to
    inter-arm handover and to a wider bimanual task set. It should be one
    deliberate pass, because the home pose, the anchor and every clearance
    figure move together.
  \item \textbf{Validate the safety layer on hardware}, starting with the
    session recovery path, which is the riskiest because the arm permits
    exactly one session.
  \item \textbf{Run the two-person study.} The exchange rate between operator
    performance and wearer trust is the open question this architecture was
    built to ask, and it remains open.
\end{enumerate}

\section{Closing}

The most transferable outcome of this work is not the mapping or the safety
architecture. It is the discipline that produced them. Four times a result
that looked like a system failure turned out to be a defect in the measuring
tool, and each was caught only by checking the instrument against something
whose answer was already known. On a robot that stands beside a person, the
difference between a measurement and an artefact is not an academic
distinction, and the habit of establishing which one is in hand before acting
on it is the part of this project most worth carrying forward.
